
seulement si "$G$ est sans torsion", i.e. $G$ est isomorphe à un produit de groupes $\operatorname{S\ell}(n)$ et $\operatorname{Sp}(m)$. Donc dans ce cas, et dans ce cas seulement, l'homomorphisme
\[
 \phi:A(G)\longrightarrow\Gtop^\bullet(G)
\]
est un isomorphisme. Donc si $G$ a de la torsion, le schéma sous-jacent fournit un exemple d'un $\phi$ non bijectif.

\textbf{4.8.} On trouve de même un exemple où l'homomorphisme $\gamma$ de (4.1)
\begin{equation}\tag{4.9}
 \gamma:A(G)\longrightarrow H^{2*}(G,\underline{Z}_\ell(*))
\end{equation}
ne se factorise pas par $\Gtop^\bullet(G)$, ou ce qui revient au même ici, un exemple où cet homomorphisme n'est pas nul en degrés $>0$. Se limitant (comme il est loisible grâce à des théorèmes de spécialisation convenables) au cas où $k=\C$, on est ramené à étudier l'homomorphisme
\begin{equation}\tag{4.10}
 A(G)\longrightarrow H^*(G(\C),\Z),
\end{equation}
déduit d'ailleurs par passage au quotient de l'homomorphisme canonique
\[
 A(G/B)\simeq H^*(G/B(\C),\Z)\longrightarrow H^*(G(\C),\Z).
\]
Il est tentant de conjecturer que l'homomorphisme (4.10) est en fait injectif, puisqu'on sait par ailleurs que les nombres premiers de torsion pour $G(\C)$ sont exactement ceux qui divisent l'ordre du groupe fini $A^+(G)=\coprod_{i>0}A^i(G)$. S'il en était ainsi, dès que $G$ aurait de la torsion, l'homomorphisme (4.10) ne se factoriserait pas par $\phi$, donc pour $\ell$ convenable (de façon précise, pour $\ell$ divisant l'ordre de $A^+(G)$, i.e. $\ell$ coefficient de torsion de $G$) il en serait de même pour (4.9). On peut du moins affirmer que dans le cas du groupe $G=\underline{G}_2$ sur le corps $\C$, l'homomorphisme (4.10) n'est pas nul en degré $>0$, de façon précise il existe un élément de $A^3(G)$ dont l'image dans $H^6(G(\C),\Z)$ est d'ordre $2$; il en résulte donc que l'homomorphisme (4.9) ne se factorise pas par $\phi$ lorsque $\ell=2$. Le fait utilisé concernant $\underline{G}_2$, qui revient à dire qu'il existe dans $H^6(G/B(\C),\Z)$ un élément dont l'image dans $H^6(G(\C),\Z)$ est non nulle, se trouve dans [3], modulo le dictionnaire connu entre le langage groupes algébriques semi-simples complexes, et groupes de Lie compacts; il m'a été communiqué obligeamment par J.P. SERRE.

